\section{Supporting evidence}\label{sec:appendix-supporting-evidence}

This appendix records the computational and audit evidence that accompanies the closed theorem package. Its role is supportive rather than foundational. None of the diagnostics below should be read as a substitute for a closed theorem statement, and none of them enlarges the claim scope of the main text.

The manuscript source, frozen audit artifacts, and build outputs are maintained in the companion repository \url{https://github.com/ioannist/six-birds-cantor}.

The supporting evidence used in the paper falls into five categories:
\begin{enumerate}
    \item the continuous pilot,
    \item the regime diagnostics,
    \item the pressure-closure checks,
    \item the strict-extension audit,
    \item and the conditional-disintegration support checks.
\end{enumerate}
All of this evidence is restricted to the audited shell. No shell-general, broader-class, or external/non-SFT breadth claim is inferred from it.

\paragraph{Continuous pilot and regime diagnostics.}
The continuous pilot and the shell-level regime diagnostics establish that the audited-shell class is not a trivial frozen regime. The manuscript-facing summaries of this evidence are Figure~\ref{fig:full-loop-regime} and Figure~\ref{fig:primitive-knockout-closure}. Figure~\ref{fig:full-loop-regime} summarizes persistent full-loop activity on the audited shell, while Figure~\ref{fig:primitive-knockout-closure} summarizes the leave-one-primitive-out degradations that certify six-mechanism causal closure. These figures support the lawfulness theoremlet, but they are not themselves theorem statements.

\paragraph{Pressure-closure checks.}
The finite-horizon shell checks that accompany the cocycle-pressure theoremlet are summarized in Table~\ref{tbl:pressure-closure-support}. Their role is to document the shell-uniform control, bounded Fekete gap proxy, and nondegenerate pressure regime used in the closed pressure package. They support the cocycle pressure closure theoremlet on the audited shell only. They do not constitute a broader-class pressure theorem.

\paragraph{Strict-extension audit.}
The audited-shell extension evidence is summarized in Table~\ref{tbl:strict-theory-extension}. The table records non-factorization, material \(P4\leftarrow P5\) forcing, packaged-strata multiplicity over \(\Tzero\)-classes, and macro-admissibility obstruction. Its purpose is to document the shell-level certificates behind the strict theory extension theoremlet. It does not support any broader-class or shell-general extension claim.

\paragraph{Conditional-disintegration support.}
The shell-level thermodynamic consequence evidence is summarized in Table~\ref{tbl:conditional-disintegration}. The table records the weighted package-conditioned pressure gap on the audited shell and shows that the gap is bounded away from zero on the two frozen witnesses. Consistent with the formulation in Section~\ref{sec:conditional-disintegration}, the evidence supports a conditional-disintegration statement, not a direct stratumwise root-separation claim.

\paragraph{Support-only closure-deficit proxy.}
The only support-only item retained in the appendix is the closure-deficit proxy table:
\begin{table}[t]
\caption{Support-only closure-deficit proxy values associated with the conditional
disintegration theoremlet.
These KL-style quantities help interpret the packaged future as
thermodynamically richer than the base cocycle object, but they are
not themselves part of any closed theorem statement in this paper.}
\label{tbl:closure-deficit-proxy}
\centering
\small
\begin{tabular}{@{}l c c@{}}
\toprule
Witness & Max.\ closure-deficit proxy & Status \\
\midrule
Base  & $3.77\times 10^{-8}$  & support-only \\
Shell & $1.08\times 10^{-11}$ & support-only \\
\bottomrule
\end{tabular}
\end{table}

This table is included because it helps interpret the packaged future as thermodynamically richer than the base cocycle object. However, it is \emph{not} part of any closed theorem statement in the paper. In particular, the manuscript does not claim an exact KL-based closure-deficit theorem.

\paragraph{Scope discipline.}
The diagnostics in this appendix are subordinate to the theoremlets of the main text and support only those four theoremlets. The non-claims collected in Section~\ref{sec:discussion} apply here unchanged.
